% \chapter*{Préface}
% \addcontentsline{toc}{chapter}{Préface}

% Il y a quatre ans, j'ai pris la décision de quitter Fribourg pour venir étudier à Yverdon-les-Bains. Ce fut un choix déterminant dans ma vie : je souhaitais sortir de ma zone de confort et apprendre à évoluer dans un environnement nouveau, moi qui avais jusqu'alors toujours eu mes écoles à cinq minutes de chez moi.

% Cette transition s'est toutefois avérée exigeante. Il m'a fallu faire preuve d'une plus grande rigueur, travailler de manière plus régulière et m'adapter à un cadre inconnu. Je me rappelle encore mon premier jour, où j'ignorais même à quel arrêt de train je devais descendre.

% Ce parcours n'a pas été sans embûches, qu'il s'agisse de trains supprimés les jours d'examens, mais surtout d'une deuxième année particulièrement éprouvante. Les résultats scolaires ne suivaient pas, et mon moral en a lourdement pâti. J'ai terminé cette année-là à bout de forces.

% L'été qui a suivi a été une période de grand vide. Pourtant, deux semaines avant la rentrée, il y a eu un déclic, un véritable besoin de changement. J'ai décidé de reprendre le contrôle et de modifier drastiquement mon approche. Ne voulant à aucun prix revivre les difficultés de ma deuxième année, je me suis imposé une routine : un travail plus régulier couplé à la pratique du sport. L'adoption de cette hygiène de vie fut sans doute la meilleure décision de mon existence, m'apportant un équilibre et un bonheur inédits.

% S'il me reste encore beaucoup de projets à concrétiser, j'ai toute la vie devant moi pour y parvenir. Ce choix d'études et d'indépendance m'a énormément appris. Malgré la rudesse de certains moments, je suis toujours là. Je suis fier du chemin parcouru et de ce que j'ai accompli. C'est cette résilience qui m'anime aujourd'hui et qui m'empêchera toujours d'abandonner.